% SYNTHETIC EVALUATION DRAFT — written for the retrieval eval in evals/.
% Near-topic distractor for peft.tex. Not a real paper; contains no results.
\documentclass{article}
\title{Outlier-Aware Post-Training Quantization and Sparsity for Efficient LLM Inference}
\begin{document}
\maketitle

\begin{abstract}
Serving large language models is bounded by memory capacity and memory
bandwidth: at batch size one, decoding a token requires reading every weight
from GPU memory. Compressing a trained model to low bit-width weights and
activations reduces both. This draft proposes a post-training quantization
method that requires no gradient updates, handles the activation outliers that
emerge in large transformers, and combines 4-bit weights with 2:4 structured
sparsity and an 8-bit key-value cache to accelerate inference.
\end{abstract}

\section{Introduction}
Large language models with tens of billions of parameters do not fit on a
single accelerator in 16-bit precision, and their autoregressive decoding is
memory-bound. Model compression lowers the cost of deployment without
retraining. Post-training quantization (PTQ) maps a trained model's weights,
and optionally its activations, to low-precision integers using only a small
calibration set, in contrast to quantization-aware training, which simulates
quantization during a full training run.

Two obstacles make PTQ of large transformers hard. First, a small number of
hidden dimensions carry activation outliers with magnitudes far larger than
the rest, so a per-tensor activation scale wastes most of the integer range.
Second, rounding each weight to its nearest grid point independently ignores
correlations between weights, and errors accumulate across layers at 3 and 4
bits.

Our contributions are: (1) an outlier-aware scheme that migrates activation
scale into the weights per channel before quantization, enabling 8-bit
activations; (2) a second-order weight quantizer that quantizes columns
sequentially and updates the remaining columns to compensate for rounding
error, using a Hessian estimated from calibration activations; and (3) joint
4-bit quantization with 2:4 semi-structured pruning, plus an 8-bit KV cache, so
that the model runs with fast integer and sparse tensor-core kernels.

\section{Related Work}
LLM.int8() keeps outlier feature dimensions in 16-bit and quantizes the rest
to 8-bit. SmoothQuant migrates quantization difficulty from activations to
weights with a per-channel scaling. GPTQ quantizes weights column by column
with an approximate second-order update derived from Optimal Brain
Quantization. AWQ protects salient weight channels identified from activation
statistics. SpQR and SqueezeLLM store sensitive weights at higher precision.
SparseGPT prunes large models in one shot with a related second-order
reconstruction. KV cache quantization reduces memory for long contexts. Our
method targets the inference-only setting and does not train any parameters.

\section{Method}
\subsection{Calibration}
We run a few hundred calibration sequences through the full-precision model
and record per-channel activation maxima and the input Gram matrix of every
linear layer.

\subsection{Outlier Migration}
For each linear layer we choose per-input-channel scales $s$ such that
$Y = (X\,\mathrm{diag}(s)^{-1})(\mathrm{diag}(s)\,W)$, balancing the dynamic
range between activations and weights. Activations are then quantized to INT8
per token and weights to INT4 per group of 128 channels.

\subsection{Error-Compensating Weight Rounding}
Columns are quantized in order of decreasing Hessian diagonal. After each
column is rounded, the residual error is distributed onto the not-yet-quantized
columns using the inverse Hessian, which keeps the layer's output close to the
full-precision output on calibration data.

\subsection{Structured Sparsity and KV Cache}
Within the same pass we enforce a 2:4 sparsity pattern by removing the two
weights with the smallest saliency in every group of four. Keys and values are
stored in INT8 with per-head scales.

\section{Experimental Plan}
We plan to compress open decoder-only language models of several sizes and
measure perplexity on held-out text, zero-shot accuracy on reasoning
benchmarks, and end-to-end decoding throughput and latency on a single GPU.
Baselines are round-to-nearest, LLM.int8(), SmoothQuant, GPTQ, AWQ and
SparseGPT at matched bit-widths and sparsity.

\section{Limitations}
Calibration data that differs from deployment traffic can bias the scales.
Speed-ups depend on kernel support for mixed INT4 and sparse formats, and
quantizing below 3 bits remains lossy for smaller models.

\end{document}
